\section*{Capítulo \ref{ch:g12:randvar} --- Variáveis aleatórias e a distribuição binomial}

\begin{solution}{exo:g12:randvar:1}
$\E(X) = -0.3 + 0 + 0.8 + 0.5 = 1$.
$\E(X^2) = 0.3 \times 1 + 0 + 0.4\times4 + 0.1\times25 = 4.4$, logo, por
König--Huygens, $\V(X) = 4.4 - 1 = 3.4$ e
$\sigma(X) = \sqrt{3.4} \approx 1.84$.
\end{solution}

\begin{solution}{exo:g12:randvar:2}
\emph{1.} As $10$ questões são \omterm{def:g12:randvar:bernoulli}{ensaios de Bernoulli} \omterm{def:g12:condprob:indep}{independentes} com
$p = \frac14$: $X \sim \mathcal B\left(10, \frac14\right)$,
$\E(X) = 2.5$,
$\sigma(X) = \sqrt{10 \times \frac14 \times \frac34} = \sqrt{1.875}
\approx 1.37$.

\emph{2.} $\P(X = 0) = \left(\frac34\right)^{10} \approx 0.056$;
\[
\P(X = 5) = \binom{10}{5}\left(\frac14\right)^5\left(\frac34\right)^5
\approx 0.058;
\quad
\P(X \geq 1) = 1 - \left(\tfrac34\right)^{10} \approx 0.944 .
\]
\end{solution}

\begin{solution}{exo:g12:randvar:3}
Ensaios idênticos e \omterm{def:g12:condprob:indep}{independentes}: $X \sim \mathcal B(400,\ 0.05)$, logo
$\E(X) = 20$ sinistros e
\[
\sigma(X) = \sqrt{400 \times 0.05 \times 0.95} = \sqrt{19} \approx 4.4 .
\]
\end{solution}

\begin{solution}{exo:g12:randvar:4}
O pagamento $P$ recebido satisfaz
$\E(P) = \frac{5 + 6}{6} = \frac{11}{6}$, de modo que o preço equitativo é
$c = \frac{11}{6} \approx 1.83$ euros. O ganho equitativo
$G = P - \frac{11}6$ assume os valores
$-\frac{11}{6}, \frac{19}{6}, \frac{25}{6}$ com \omterm{def:g10:proba:distribution}{probabilidades}
$\frac46, \frac16, \frac16$:
\[
\V(G) = \E(G^2)
= \frac{4 \times 121 + 361 + 625}{6 \times 36} = \frac{1470}{216}
= \frac{245}{36} \approx 6.8,
\qquad \sigma(G) \approx 2.6 .
\]
Um jogo equitativo tem ganho \emph{médio} nulo, mas seus resultados ainda
flutuam: equitativo não é o mesmo que sem risco.
\end{solution}

\begin{solution}{exo:g12:randvar:5}
$X \sim \mathcal B(8,\ 0.7)$.
\[
\P(X = 6) = \binom86 (0.7)^6(0.3)^2 \approx 0.296 ;
\]
\[
\P(X \geq 6) = \P(6) + \P(7) + \P(8)
\approx 0.296 + 8(0.7)^7(0.3) + (0.7)^8
\approx 0.296 + 0.198 + 0.058 = 0.552 ;
\]
$\P(X \geq 1) = 1 - (0.3)^8 \approx 0.99993$. Moda:
$(n+1)p = 6.3$, de modo que o valor mais provável é $k^* = 6$
(\cref{exo:g12:randvar:7}).
\end{solution}

\begin{solution}{exo:g12:randvar:6}
O número de passageiros que comparecem é $X \sim \mathcal B(104,\ 0.9)$; o
voo fica com excesso de passageiros quando $X \geq 101$:
\[
\P(X \geq 101) = \sum_{k=101}^{104}\binom{104}{k}(0.9)^k(0.1)^{104-k}
\approx 0.006 .
\]
Vender $4\%$ mais passagens do que assentos provoca um incidente em apenas
cerca de $0.6\%$ dos voos --- a economia por trás do excesso de reservas.
\end{solution}

\begin{solution}{exo:g12:randvar:7}
\[
\frac{\P(X = k+1)}{\P(X = k)}
= \frac{\binom{n}{k+1}}{\binom nk}\cdot\frac{p}{1-p}
= \frac{n - k}{k + 1}\cdot\frac{p}{1-p},
\]
usando $\binom{n}{k+1} = \binom nk \frac{n-k}{k+1}$. Essa razão é
$\geq 1$ se, e somente se, $(n-k)p \geq (k+1)(1-p)$, se, e somente se,
$np - k p \geq k - kp + 1 - p$, se, e somente se, $k \leq (n+1)p - 1$.
Assim, as \omterm{def:g10:proba:distribution}{probabilidades} crescem estritamente enquanto
$k + 1 \leq (n+1)p$ e decrescem depois: o \omterm{def:g10:functions:extrema}{máximo} é atingido em
$k^* = \floor{(n+1)p}$ (compartilhado com $k^* - 1$ quando $(n+1)p$ é
inteiro).
\end{solution}

\begin{solution}{exo:g12:randvar:8}
\emph{1.} A primeira cara no lançamento $k$ significa $k-1$ coroas
seguidas de uma cara: \omterm{def:g10:proba:distribution}{probabilidade}
$\left(\frac12\right)^{k-1}\cdot\frac12 = 2^{-k}$, para
$1 \leq k \leq 10$. \omterm{def:g10:proba:distribution}{Probabilidade} total de ganhar:
$\sum_{k=1}^{10} 2^{-k} = 1 - 2^{-10} = \frac{1023}{1024}$ (o jogo só é
perdido com dez coroas consecutivas).

\emph{2.} Ganho esperado:
\[
\sum_{k=1}^{10} 2^k \cdot 2^{-k} = \sum_{k=1}^{10} 1 = 10 \text{ euros}.
\]
Sem o limite, a soma $\sum_{k\geq1} 1$ diverge: o ganho esperado é
infinito, embora o jogo quase sempre pague uma quantia pequena --- o
famoso \emph{paradoxo de São Petersburgo}, que mostra que a \omterm{def:g12:randvar:exp}{esperança}
sozinha não mede o valor de um jogo.
\end{solution}

\begin{solution}{pb:g12:randvar:1}
\textbf{1.} $X \sim \mathcal B\!\left(3, \frac16\right)$:
$\P(X = 0) = \frac{125}{216}$,
$\P(X = 1) = \frac{75}{216}$,
$\P(X = 2) = \frac{15}{216}$,
$\P(X = 3) = \frac{1}{216}$.

\textbf{2.} $\E(X) = 3 \times \frac16 = \frac12$;
$V(X) = 3 \times \frac16 \times \frac56 = \frac{5}{12}$.

\textbf{3.} Pagamento esperado $2\,\E(X) = 1$ euro contra uma
aposta de $1$ euro: ganho $0$ --- exatamente equitativo (uma
raridade).

\textbf{4.} Sem reposição, as retiradas são dependentes (as
chances da segunda carta dependem da primeira): não é binomial
--- falha o item da \omterm{def:g12:condprob:indep}{independência} na lista de verificação. Com
reposição: $n = 5$ fixo, $p = \frac14$ constante, retiradas
\omterm{def:g12:condprob:indep}{independentes}: binomial.

\textbf{5.} $\E(X) = 6$, $\sigma(X) = \sqrt{4.2} \approx
2.05$. $\E(S) = 5 \times 6 - 10 = 20$;
$\sigma(S) = 5\sigma(X) \approx 10.25$.

\textbf{6.} $n = 105$ ensaios fixos (as passagens), cada um um
comparecimento ou não com o mesmo $p = 0.9$, supostos
\omterm{def:g12:condprob:indep}{independentes}: binomial. A confissão: as famílias perdem voos
\emph{juntas} e as tempestades esvaziam aviões inteiros --- a
\omterm{def:g12:condprob:indep}{independência} é o ato de fé do modelo, e ausências
correlacionadas engordam as caudas mais do que a binomial
promete.

\textbf{7.} $\E(X) = 94.5$;
$\sigma(X) = \sqrt{105 \times 0.9 \times 0.1} =
\sqrt{9.45} \approx 3.07$.

\textbf{8.} $\P(X \geq 101) = \sum_{k=101}^{105}
\binom{105}{k} 0.9^k\, 0.1^{105 - k}$.

\textbf{9.} $\approx 0.0167$: cerca de um voo em sessenta
preteriu alguém.

\textbf{10.} $\E(\text{preteridos}) \approx 0.023$ passageiros
por voo: indenização esperada $\approx 19$ euros --- contra
$1\,000$ euros de receita extra. Vender cinco passagens a mais
é esmagadoramente lucrativo; daí toda companhia aérea fazer
isso.

\textbf{11.} Pela tabela de \omterm{def:g10:proba:distribution}{probabilidades} de cauda:
$n = 105$: $1.7\,\%$; $n = 106$: $4.0\,\%$; $n = 107$:
$8.1\,\%$. A maior venda em conformidade é de $n = 106$
passagens.

\textbf{12.} $z = \frac{100.5 - 94.5}{3.07} \approx 1.95$: o
limiar fica dois \omterm{def:g11:stat:variance}{desvios padrão} acima da \omterm{def:g11:stat:mean}{média}, e a regra
prática da curva em sino (``cauda superior de $2\sigma$
$\approx 2.5\,\%$'') chega perto do valor exato de $1.7\,\%$
--- a curva suave que sombreia a binomial é a protagonista dos
próximos capítulos.

\textbf{13.} $0.98^{20} + 20 \times 0.02 \times 0.98^{19}
\approx 0.94$: o lote honesto passa em $94\,\%$ das vezes.

\textbf{14.} $0.9^{20} + 20 \times 0.1 \times 0.9^{19}
\approx 0.39$: o lote ruim ainda escapa em $39\,\%$ das vezes.

\textbf{15.} Risco do produtor: um lote bom rejeitado
($\approx 6\,\%$); risco do consumidor: um lote ruim aceito
($\approx 39\,\%$). Uma \omterm{def:g10:proba:sample}{amostra} maior (com limiar de aceitação
proporcional) reduz os dois --- ao custo de mais inspeção: a
qualidade tem uma linha de orçamento.

\textbf{16.} $V\!\left(\frac Xn\right) = \frac{V(X)}{n^2} =
\frac{np(1-p)}{n^2} = \frac{p(1-p)}{n}$: \omterm{def:g11:stat:variance}{desvio padrão}
$\sqrt{\frac{p(1-p)}{n}}$. Quadruplique a \omterm{def:g10:proba:sample}{amostra}, reduza o
ruído à metade: a lei do $\frac{1}{\sqrt n}$ dos \omterm{def:g10:numbers:interval}{intervalos} de
flutuação do ano 10, agora demonstrada.

\textbf{17.} O limite prende o pagamento em $2^{10} = 1024$
euros, de modo que cada uma das dez rodadas contribui com
exatamente $1$ euro de \omterm{def:g12:randvar:exp}{esperança}: $\E = 10$. A \omterm{def:g12:randvar:exp}{esperança}
infinita do jogo sem limite vinha inteiramente de ganhos
astronomicamente raros e astronomicamente grandes; pôr um
limite é confessar que nenhum banco paga $2^{50}$ euros. Um
preço equitativo do bilhete para o jogo com limite: $10$ euros
--- e, de repente, ninguém mais fica paradoxado.

\textbf{18.} Prêmio esperado por bilhete:
$\frac{100 + 2 \times 50}{200} = 1$ euro contra um bilhete de
$2$ euros: margem de $50\,\%$ --- dezoito vezes os $2.7\,\%$
da roleta. A rifa sobrevive porque seus participantes estão
doando conscientemente: a ``perda'' é o objetivo.

\textbf{19.} Sinistros esperados: $1\,000 \times 0.01 \times
10\,000 = 100\,000$; prêmios $130\,000$: lucro esperado
$30\,000$ euros. \omterm{def:g11:stat:variance}{Desvio padrão} do total de sinistros:
$10\,000 \times \sqrt{1\,000 \times 0.01 \times 0.99}
\approx 31\,500$ euros --- um único ano ruim comum devora todo
o lucro esperado. Daí as reservas de capital, o resseguro e
carteiras muito maiores que mil apólices: as seguradoras vivem
da lei dos grandes números e guardam reservas contra a lentidão
dela.

\textbf{20.} Lista de verificação primeiro --- $n$ fixo, mesmo
$p$, \omterm{def:g12:condprob:indep}{independência} \emph{alegada} --- ou não há binomial
nenhuma. Depois navegue por $\E \pm \sigma$: as \omterm{def:g11:stat:mean}{médias}
localizam, os desvios advertem. Depois precifique as caudas:
preterições, lotes ruins e anos ruinosos vivem todos além de
$2\sigma$. Ressalvas: a \omterm{def:g12:condprob:indep}{independência} é a promessa de quem
modela, não do mundo; e a \omterm{def:g12:randvar:exp}{esperança} ignora exatamente o que
destrói você. As páginas que faltam no regulamento --- com que
rapidez as \omterm{def:g10:stats:series}{frequências} se assentam e que forma tomam as
flutuações --- são os dois próximos capítulos.
\end{solution}
